\scanpage{89}
\mspage{6}
% archive sheet number: 88
\[
{\unsure{\Phi|\varphi}}^i_\nu=\sum_{t\in T}(\underline\Phi^i_\nu)_t,\quad\deleted{\text{avec}}
\tag*{(1.18)}
\]
où
\[
(\underline\Phi^i_\nu)_t=\varinjlim_\alpha\underline H^i_{X_{\alpha0}}(\Lambda_{\nu X_\alpha})_{t_\alpha}
\tag*{(1.19)}
\]
où $t_\alpha$ est l'unique point de $X_\alpha$ au dessus de $t$.

On voit donc que dans le cas où l'ensemble $T$ des pts de $X_0$ en lesquels $f$ est non lisse est \uline{fini}, le \deleted{calcul \unsure{du|des}} groupe $\Phi^i_\nu$ des cycles évanescents est \inserted{une somme de termes} \uline{purement locaux} \deleted{\unsure{en les points de}}\inserted{(au sens de la top. étale)} relatifs aux différents points de $\unsure{t|T}$). \deleted{\unsure{1.19}} Nous donnerons plus tard des exemples \deleted{de calculs} d'application de ce principe à des cas particuliers intéressants.

\inserted{1.20.} Mais nous allons nous tourner \inserted{d'abord} vers une \deleted{autre}\inserted{variante} \deleted{méthode} de la méthode précédente pour \deleted{façon de} formuler les relations entre la cohomologie de \deleted{$\unsure{\mathcal E|E}$} $X_s$ et de $X_{\bar{\unsure{\eta|y}}}$, \inserted{qui sera} plus commode dans la démonstration du th. de monodromie.

Notons toutefois \inserted{(\unsure{tout de suite})} que \inserted{On constate aisément que \unsure{.|:|;}} par construction même, le \deleted{\unsure{2.}} groupe d'inertie $I=\operatorname{Gal}(\bar K/K)$ opère sur les groupes $\Phi^i_\nu$ et \deleted{\unsure{les}} sur les faisceaux $\underline\Phi^i_\nu$, que la suite spectrale (1.14) peut être interprétée comme une suite spectrale de $I$-modules, et la suite exacte (1.9) comme une suite exacte de $I$-modules, étant entendu que $I$ opère trivialement sur $H^*(X_s,\Lambda_\nu)$. \inserted{De \unsure{même} après passage à la limite.} Il sera possible de \unsure{prouver} le th. de
